\section{The articulation measure: charge derived from articulation status}

\texttt{Language.Spine} names a gap directly: nothing in that module derives
either coordinate of a charge, and any reading of a supplied \texttt{Charge}
or \texttt{Record} as evidence remains a caller-side requirement. This
section closes that gap for the one coordinate that needs no caller at all:
an articulation cell's own status record already \emph{is} an evidence pair,
once read through the corner convention (\cref{thm:corner}).

\begin{proposition}[The measured charge]\label{thm:measure}\lean{Obsmetrologia.Articulating.Cell.measure}\leanok\uses{thm:articulation-cell-corner, thm:charge}
\texttt{measure}~\(c\) is the charge \(c.\mathrm{corner}\): the corner a cell
already sits in, handed back as its charge. Since \(\mathtt{Charge} :=
\mathtt{Articulating.Corner}\) definitionally, this is the identity read at
two depths --- trivial as a construction, but name-load-bearing, because it
lets a caller obtain a charge \emph{from the articulation record itself},
where \texttt{Spine.posit} otherwise demands one admitted from outside.
\end{proposition}

\begin{proposition}[The measured meme]\label{thm:meme}\lean{Obsmetrologia.Articulating.Cell.meme}\leanok\uses{thm:measure}
\texttt{meme}~\(c\) is the posit at \(c\)'s own claim
(\(\texttt{Claim.seed}\,c.\mathrm{utterance}\)) whose charge is
\texttt{measure}~\(c\) --- derived from the cell's own status, not admitted
by a caller, unlike \texttt{Spine.posit}.
\end{proposition}

\begin{proposition}[The four tellings, by name]\label{thm:measure-tellings}\lean{Obsmetrologia.Articulating.Cell.measure_absurd, Obsmetrologia.Articulating.Cell.measure_potential, Obsmetrologia.Articulating.Cell.measure_active, Obsmetrologia.Articulating.Cell.measure_stable}\leanok\uses{thm:measure, thm:four-charges}
Measuring the four corners recovers the four distinguished charges by name.
The un-articulated and articulated corners need no witness beyond the
utterance itself: \(\texttt{measure}\,(\texttt{unarticulated}\,u) = \bot_k\)
and \(\texttt{measure}\,(\texttt{articulate}\,u) = t\). The active and stable
corners are read off an explicit token witness \(w\), since neither is what
articulating alone produces: \(\texttt{measure}\,\langle u, [], [w]\rangle =
f\) --- the active corner, exhibited by fiat as a value-algebra construction,
not an articulation --- and \(\texttt{measure}\,(\texttt{engage}\,
(\texttt{articulate}\,u)\,w) = \top_k\).
\end{proposition}

\begin{theorem}[The un-articulated cell measures --- and posits --- to bottom]\label{thm:meme-initial-charge}\lean{Obsmetrologia.Articulating.Cell.meme_initial_reading}\leanok\uses{thm:meme, thm:measure-tellings, thm:articulation-cell-initial, thm:kBot-least}
The un-articulated cell (\texttt{unarticulated}~\(u\), born absurd, told
none --- \cref{thm:articulation-cell-initial}, never articulating's own
output) measures --- and posits --- to \(\bot_k\): the derived reading and
the knowledge-least reading of \cref{thm:knowledge-least} agree at the one
point both are pinned.
\end{theorem}

\begin{definition}[The status order on cells]\label{thm:statusle}\lean{Obsmetrologia.Articulating.Cell.statusLe}\leanok\uses{thm:articulation-cell}
\texttt{statusLe}~\(c\)~\(d\) holds when each of \(c\)'s two articulation
statuses implies the like status of \(d\). It compares evidence tellings
directly on the status pairs; no shared utterance is required, since this is
not a comparison of two cells over one claim.
\end{definition}

\begin{theorem}[To articulate is to raise the measured charge]\label{thm:measure-mono}\lean{Obsmetrologia.Articulating.Cell.measure_mono}\leanok\uses{thm:statusle, thm:measure, thm:kle}
If \(d\) tells everything \(c\) tells (and possibly more), \(d\)'s measured
charge sits above \(c\)'s in the knowledge order: \(\texttt{statusLe}\,c\,d
\to \bot_k \le_k \dots\), exhausted over the four \((\mathrm{token},
\mathrm{claim})\) bools of each cell.
\end{theorem}

\begin{definition}[The delta valuation of a cell]\label{thm:deltavaluation}\lean{Obsmetrologia.Articulating.Cell.deltaValuation}\leanok\uses{thm:measure, thm:knowledge-least}
The knowledge-least valuation, except at \(c\)'s own claim, where it carries
\(c\)'s measured charge: the exact shape the meme-folding valuation trace
advances by one step, specialized to a knowledge-least base.
\end{definition}

\begin{theorem}[Articulation history is visible one level up]\label{thm:deltavaluation-mono}\lean{Obsmetrologia.Articulating.Cell.deltaValuation_mono}\leanok\uses{thm:deltavaluation, thm:measure-mono}
At a \emph{shared} claim --- \(c\) and \(d\) over the same utterance --- a
status-increasing pair of cells measures to a \(\le_k\)-increasing pair of
delta valuations: articulation history driving one claim's charge upward is
visible one level up, as a pointwise \(\le_k\) step between valuations.
\end{theorem}

\begin{theorem}[The measured-vs-seeded semantics question]\label{thm:measured-vs-seeded}\notready\uses{thm:meme-initial-charge, thm:knowledge-least}
This section derives one coordinate --- the charge a cell's own articulation
status supplies --- from that status alone. It does not resolve, and is not
built to resolve, whether a \emph{measured} charge (\cref{thm:measure}) and
an independently \emph{seeded}, caller-supplied charge (\texttt{Spine.posit})
at the same claim should ever be required to agree, or how the two interact
under settling. \texttt{docs/history/corner-naming-convention.md} records
this exact question as open --- the closure-of-assertions reading against
the LFP-of-the-measured-carrier reading --- and this section leaves it there,
open, rather than picking a side.
\end{theorem}
